\documentclass[11pt,a4paper,landscape]{article}

\usepackage{fontawesome5}
\providecommand{\faPencilAlt}{\faIcon{pencil-alt}}

\usepackage{styles/couleurs}
\usepackage{styles/mise_en_page}
\usepackage{styles/boites}
\usepackage{styles/maths}
\usepackage{styles/figures}

\usepackage{enumitem}
\usepackage{array}
\usepackage{tabularx}
\usepackage{tikz}
\usepackage{graphicx}
\usepackage{amssymb}
\usepackage{multicol}

\usetikzlibrary{calc}

%==================================================
% Réglages : A4 paysage
%==================================================

\geometry{
  a4paper,
  landscape,
  top=0.8cm,
  bottom=0.8cm,
  left=1cm,
  right=1cm
}

\pagestyle{empty}

\setlength{\parindent}{0pt}
\setlength{\parskip}{0.12em}
\setlength{\columnsep}{0.85cm}

\setlist[enumerate]{
  leftmargin=*,
  itemsep=0.05em,
  topsep=0.08em,
  parsep=0pt
}

\titlespacing*{\section}
  {0pt}
  {0.35em}
  {0.12em}

%==================================================
% Petits outils de mise en page
%==================================================

\newcommand{\soustitreexo}[1]{%
  {\sffamily\bfseries\color{TitreC}#1}\par
  \vspace{0.04em}
}

\newcommand{\reponse}[1][2.4cm]{%
  \rule{#1}{0.35pt}
}

\newcommand{\lignepointillee}{%
  \noindent\dotfill\par
  \vspace{0.18em}
}

\newcommand{\deuxlignes}{%
  \lignepointillee
  \lignepointillee
}

%==================================================
% Document
%==================================================

\begin{document}

\begin{center}
  {\sffamily\LARGE\bfseries\color{TitreC}
  Séance d'AP -- Statistiques}\\[-0.05em]

  {\sffamily\normalsize\bfseries\color{SavoirC}
  Vocabulaire, effectifs et moyenne}
\end{center}

\vspace{-0.45em}

\begin{multicols*}{2}

%==================================================
% 1. VOCABULAIRE ET EFFECTIFS
%==================================================

\section*{1. Vocabulaire et effectifs}

\soustitreexo{Exercice 1}

On a demandé à \(24\) élèves combien de livres ils avaient lus
pendant les vacances :

\[
0;\ 2;\ 1;\ 0;\ 3;\ 1;\ 2;\ 0;\ 1;\ 4;\ 2;\ 1;\
0;\ 3;\ 1;\ 2;\ 5;\ 1;\ 2;\ 0;\ 1;\ 3;\ 2;\ 1
\]

\begin{enumerate}[label=\alph*)]

  \item Population :
  \reponse[4.2cm]

  \item Caractère étudié :
  \reponse[3.4cm]

  \item Valeurs prises par le caractère :
  \reponse[3.1cm]

  \item Effectif de la valeur \(2\) :
  \reponse[1.8cm]

\end{enumerate}


\soustitreexo{Exercice 2}

On a interrogé des élèves sur leur mode de transport habituel
pour venir au collège.

\begin{center}

\renewcommand{\arraystretch}{1.12}

\begin{tabular}{|c|c|c|c|c|}
\hline

Mode de transport
& À pied
& Bus
& Vélo
& Voiture
\\
\hline

Effectif
& 6
& 11
& 4
& 7
\\
\hline

\end{tabular}

\end{center}

\begin{enumerate}[label=\alph*)]

  \item Effectif des élèves venant en bus :
  \reponse[1.7cm]

  \item Effectif total :
  \reponse[1.8cm]

  \item Nombre d'élèves ne venant pas en bus :
  \reponse[1.7cm]

\end{enumerate}


%==================================================
% 2. LECTURE DE DOCUMENT
%==================================================

\section*{2. Lire un document}

\soustitreexo{Exercice 3}

Le diagramme indique le nombre d'élèves inscrits
à quatre activités.

\begin{center}

\begin{tikzpicture}[x=1.35cm,y=0.27cm]

  \draw[->,line width=0.7pt]
    (0,0) -- (4.7,0);

  \draw[->,line width=0.7pt]
    (0,0) -- (0,10.5)
    node[above] {\small Effectif};

  \foreach \y in {0,2,...,10}{
    \draw[NeutreC!18]
      (0,\y) -- (4.35,\y);

    \node[left,font=\small]
      at (0,\y) {\y};
  }

  \fill[SavoirC!70]
    (0.45,0) rectangle (1.05,8);

  \fill[SavoirC!70]
    (1.45,0) rectangle (2.05,5);

  \fill[SavoirC!70]
    (2.45,0) rectangle (3.05,9);

  \fill[SavoirC!70]
    (3.45,0) rectangle (4.05,6);

  \node[below,font=\small]
    at (0.75,0) {Foot};

  \node[below,font=\small]
    at (1.75,0) {Danse};

  \node[below,font=\small]
    at (2.75,0) {Natation};

  \node[below,font=\small]
    at (3.75,0) {Basket};

\end{tikzpicture}

\end{center}

\begin{enumerate}[label=\alph*)]

  \item Effectif en natation :
  \reponse[1.5cm]

  \item Activité d'effectif le plus faible :
  \reponse[2.4cm]

  \item Effectif total représenté :
  \reponse[1.5cm]

\end{enumerate}


%==================================================
% 3. MOYENNE
%==================================================

\section*{3. Moyenne}

\soustitreexo{Exercice 4 -- Une petite série}

Charles a obtenu les notes suivantes :

\[
9;\ 5;\ 5;\ 9;\ 9;\ 9;\ 5;\ 6;\ 2;\ 6
\]

\textbf{Calculer la moyenne puis interpréter le résultat.}

\deuxlignes


\soustitreexo{Exercice 5 -- Une série plus longue}

En octobre 2002, à Rome, on a relevé les températures suivantes :

\[
16;\ 14;\ 15;\ 16;\ 16;\ 17;\ 15;\ 14;\ 12;\ 14;\ 14;\ 14;\
15;\ 13;\ 11;\ 11
\]

\[
10;\ 9;\ 11;\ 11;\ 11;\ 13;\ 14;\ 16;\ 16;\ 15;\ 13;\
14;\ 14;\ 14;\ 16
\]

\textbf{Calculer la température moyenne puis interpréter le résultat.}

\deuxlignes


\soustitreexo{Exercice 6 -- Des données regroupées}

Pour passer une commande de chaussures de foot,
Elsa a noté les pointures des membres de son club :

\begin{center}

\renewcommand{\arraystretch}{1.15}

\begin{tabular}{|c|c|c|c|c|c|}
\hline

Pointure
& 35
& 36
& 37
& 38
& 39
\\
\hline

Effectif
& 7
& 9
& 5
& 2
& 7
\\
\hline

\end{tabular}

\end{center}

\textbf{Calculer la pointure moyenne puis interpréter le résultat.}

\deuxlignes


%==================================================
% 4. FRÉQUENCES
%==================================================



\end{multicols*}


%==================================================
% CORRECTION RAPIDE PROFESSEUR -- NON IMPRIMÉE
%==================================================
%
% EXERCICE 1
%
% Population : les 24 élèves interrogés.
%
% Caractère : nombre de livres lus pendant les vacances.
%
% Valeurs : 0 ; 1 ; 2 ; 3 ; 4 ; 5.
%
% Effectif de 2 : 6.
%
%
% EXERCICE 2
%
% Bus : 11.
%
% Effectif total :
% 6 + 11 + 4 + 7 = 28.
%
% Hors bus :
% 28 - 11 = 17.
%
%
% EXERCICE 3
%
% Natation : 9.
%
% Plus faible effectif : danse, 5.
%
% Effectif total :
% 8 + 5 + 9 + 6 = 28.
%
%
% EXERCICE 4
%
% Somme :
% 65.
%
% Moyenne :
% 65 / 10 = 6,5.
%
% Interprétation possible :
% si les 65 points étaient répartis également
% sur les dix notes, chaque note vaudrait 6,5/20.
%
%
% EXERCICE 5
%
% Somme :
% 424.
%
% Moyenne :
% 424 / 31 ≈ 13,7 °C.
%
% Interprétation :
% durant le mois, la température quotidienne
% était en moyenne d'environ 13,7 °C.
%
%
% EXERCICE 6
%
% Somme pondérée :
%
% 35×7 + 36×9 + 37×5 + 38×2 + 39×7
%
% = 1103.
%
% Effectif total :
%
% 7 + 9 + 5 + 2 + 7 = 30.
%
% Moyenne :
%
% 1103 / 30 ≈ 36,8.
%
% Interprétation :
% la pointure moyenne des membres du club
% est d'environ 36,8.
%
% Une moyenne n'est pas nécessairement
% une valeur observée dans la série.
%
%
% EXERCICE 7
%
% Foot :
%
% 10 / 25
%
% = 2/5
%
% = 0,4
%
% = 40 %.
%
% Interprétation :
% 40 % des élèves du groupe pratiquent
% principalement le football.
%

\end{document}